% Preface: the README's About, During the Course, Learning Outcomes, Written Papers and No Hardware
% Yet sections, and the course plan and software list from info/README.md, as the front matter of
% a book.

\chapter{Preface}
\markboth{Preface}{Preface}

This book is for junior embedded developers, and for developers with some programming experience
who want to learn embedded systems programming. It teaches embedded development in C on one real
target, the ATmega328P, the microcontroller on the Arduino Uno and Nano, and it goes from the C
language itself, through registers, interrupts and interrupt-driven peripheral drivers, to
hand-written encapsulation and polymorphism.

The topics, in the order the book takes them:
\begin{itemize}
  \item A C fundamentals refresher: structs, pointers, function pointers, dynamic memory, macros.
  \item Register-level I/O: bitwise operations, \code{DDRx}/\code{PORTx}/\code{PINx}.
  \item Interrupts: external interrupts and pin change interrupts.
  \item Analog-to-digital conversion (ADC).
  \item Timers: prescalers, Normal and CTC modes, non-blocking timing.
  \item Serial communication (UART).
  \item EEPROM: persisting state across resets.
  \item The watchdog timer.
  \item Structs as drivers: bundling state for peripherals with multiple instances.
  \item Encapsulation and polymorphism in C: opaque structs, function pointers in structs, and
    hand-written vtables.
\end{itemize}

Unlike a general C text, every chapter builds toward that one target, and every worked example is a
real program, compiled and flashed to actual hardware rather than simulated. The book ends where
\emph{Modern Embedded C++}, a separate follow-on course and book, begins.

\textbf{Along the way you will build circuits and write programs for them, not just read about
them.} You will:
\begin{itemize}
  \item Wire and program real circuits (LEDs, buttons, potentiometers) on an Arduino Uno or Nano.
  \item Write drivers directly against the ATmega328P's registers and datasheet.
  \item Move from polling to interrupt-driven, non-blocking designs.
  \item Bundle a peripheral's state into a struct to support multiple instances.
  \item Refactor a driver from a plain struct into an opaque type, then into a hand-written
    interface with multiple implementations.
\end{itemize}
The examples and exercises stay close to typical embedded problems:
\begin{itemize}
  \item Digital I/O and debounced button input.
  \item Timer-driven, non-blocking LED patterns.
  \item Serial reporting of sensor readings.
  \item Persisting state to EEPROM across resets.
  \item A GPIO driver evolving from a public struct, to an opaque type, to a vtable-based
    interface.
\end{itemize}

\section*{What you should be able to do at the end}
Having worked through the book, you should be able to:
\begin{itemize}
  \item Write and debug C code directly against a microcontroller's registers and datasheet.
  \item Move from blocking, polling-based designs to interrupt-driven, non-blocking ones.
  \item Keep interrupt handlers safe by flagging events in the ISR and handling them in
    \code{main()}.
  \item Bundle a peripheral's state into a struct to support several instances of the same
    peripheral.
  \item Split a driver across a header and source file, and use opacity to protect its invariants.
  \item Store and call through function pointers, including as struct members.
  \item Understand, and hand-implement, the vtable mechanism C++ uses for virtual dispatch.
  \item Move directly on to \emph{Modern Embedded C++}.
\end{itemize}

\section*{What you are assumed to know}
You are expected to be able to program already, in C or in another language. Chapter~1 is a
fast-moving refresher on C for programmers, not a first course in programming: it covers the syntax
quickly and slows down only where C differs from the languages you may know, above all at pointers.
Nothing about embedded systems or electronics is assumed beyond knowing what an LED, a button and a
resistor are. Everything else, from what a register is onwards, is taught from the beginning.

\section*{How the book is organised}
The book is typeset from a ten-lecture course, and each chapter is one lecture. Its sections are the
lecture's appendices, with the same numbered subsections, so a chapter here and a lecture directory
in the repository can be read side by side: the lecture's Appendix~A.7 is a chapter's first section,
subsection~7.

\begin{booktable}{@{}p{5.2em}p{3.6em}L@{}}
  \thead{Chapter} & \thead{Lecture} & \thead{Topic} \\
  \midrule
  1 & \file{L01} & C fundamentals refresher \\
  2 & \file{L02} & Registers and I/O \\
  3 & \file{L03} & Interrupts \\
  4 & \file{L04} & Analog-to-digital conversion (ADC) \\
  5 & \file{L05} & Timers \\
  6 & \file{L06} & Serial communication (UART) \\
  7 & \file{L07} & EEPROM \\
  8 & \file{L08} & The watchdog timer \\
  9 & \file{L09} & Structs as drivers \\
  10 & \file{L10} & Encapsulation and polymorphism in C \\
\end{booktable}

Chapters~1 to~6 are the introductory arc: the language, digital I/O, reacting to events, reading
the analog world, timing without blocking, and reporting back to a PC. Chapters~7 and~8 complete
the survey of peripherals with the two the chip has exactly one of, and Chapters~9 and~10 are about
driver design, built on everything before them.

Every chapter opens with what is in it, works through the material, closes with a review of what
you should now be able to explain, and ends with its exercises. The appendices hold the guide to
running the demos in a simulator, the answers to the exercises that have one right answer, and the
two written papers described below.

\section*{What the exercises look like}
The exercises come in \textbf{sets}, grouped the way the lecture groups them: questions about the
concepts first, then programs. Every exercise is labelled with its kind:

\begin{booktable}{@{}p{7.4em}L@{}}
  \thead{Kind} & \thead{What it asks of you} \\
  \midrule
  \kindtag[fillaccent2]{Reflection} & Answer in prose: explain what the code does, why it is
    written the way it is, or why one design is better than another. \\
  \kindtag[fillaccent3]{Calculation} & Work a number out on paper: a voltage, a tick period, a
    register value. \\
  \kindtag{Code} & Write a program, build it, and run it on the board, or in the simulator where
    there is no board. \\
\end{booktable}

\textbf{The Reflection and Calculation exercises are answered in Appendix~\ref{app:answers}},
chapter by chapter, calculations with their working. \textbf{The Code exercises have no printed
solution.} Each is a program that either does what the exercise describes on the board or does not,
and the board is how you check it. Some Code exercises also ask a question about the design they
make you write; the appendix answers those questions, but not with code. Try each exercise before
you read its answer.

\section*{Two written papers, and what they are not}
Nothing in the course this book comes from is marked. Assessment is the exercise set after every
chapter, and the questions in every chapter's review.

Appendices~\ref{exam:ch:about} to \ref{ansb:ch:answers} hold two three-hour papers with worked
solutions, and they check something else: \textbf{what you can reconstruct on paper, with nothing
in front of you.} Eight questions each, mixing theory with C you either read and repair or write
from scratch, and no compiler, datasheet or board in the room to tell you which it is.

\textbf{They exist purely so that you can test your own knowledge after working through the book.
They gate nothing, they are not a qualification, and no part of the book requires them.}

\textbf{Take one after the last chapter.} Both papers draw on all ten chapters, so sitting one
partway through examines material you have not read yet.

\section*{What you need}
Every worked example in this book runs on one board, and you need one to follow it the way it is
written:
\begin{itemize}
  \item \textbf{An Arduino Uno or Nano.} Either works: the same ATmega328P, the same registers, the
    same pin numbering. \Secref{c2:a:2} compares the two.
  \item \textbf{A breadboard and a handful of parts}: LEDs with 220--470\,\textohm{} series
    resistors, push buttons with 10\,k\textohm{} pull-down resistors, a potentiometer for
    Chapter~4, and jumper wires. Chapter~6's last exercise adds a TMP36 temperature sensor.
\end{itemize}
And the software to build and flash programs for it, on either of two paths:
\begin{itemize}
  \item \textbf{On Windows, Microchip Studio}, the IDE for writing and building the examples, with
    the Arduino IDE used once to borrow the command that flashes the board. \Secref{c2:app:c} walks
    through it.
  \item \textbf{On Linux, or on Windows through WSL}, the AVR toolchain itself: \code{avr-gcc},
    \code{avr-objcopy} and \code{avrdude}, used from the command line with a Makefile.
    \Secref{c2:app:d} walks through it.
\end{itemize}
The companion repository holds the lectures, the worked demos and everything the book is typeset
from. Clone it:

\begin{shell}
git clone https://github.com/qrtech-academy/embedded-c.git
\end{shell}

\noindent Two commands at its root cover everything the book asks of it on Linux:

\begin{shell}
make build               # Build every demo, one project directory at a time.
make format-check        # Check the formatting of every C file.
\end{shell}

\noindent Each demo directory also builds on its own, and flashes the board with \code{make flash};
Chapter~6 shows the layout they all share.

\section*{No hardware yet?}
Every worked demo can also be run on your own computer, under the simavr simulator, including its
serial output, button presses, EEPROM that survives a restart, and debugging with \code{avr-gdb}.
Appendix~\ref{sim:ch} shows how. It is a way to keep moving until you have a board, not a
replacement for one: a simulator will not teach you what a floating input, a missing pull-up or a
badly seated jumper wire feels like. Get a board when you can.
